\section{Conclusion and Discussion}
\label{sec:discussion}

An automated architecture search surfaced forking: a replay-aligned divergence between training and validation loss in models with over-encoding $n$-gram memory. Frequency-resolved measurements tie the gap to incomplete continuation coverage of low-frequency contexts, Good--Turing missing mass predicts its frequency dependence, and replay-driven logit sharpening explains how repeated reinforcement converts coverage gaps into validation damage. The effect reproduces in a DeepSeek-style Engram model (Appendix~\ref{app:deepseek}) and reopens during continued training and verifier-selected post-training on small replayed corpora (Section~\ref{sec:transfer}).

Two practical consequences follow. First, local encoding capacity should be evaluated jointly with continuation coverage and data reuse: a smaller training--validation gap alone is not evidence of better generalization, and frequency-resolved probes expose when local gains come at the expense of uncovered continuations. Second, post-training stages that replay narrow corpora inherit the risk even from a nearly gap-free checkpoint; freezing the local memory mitigates but does not remove it, so replay schedules and gap trajectories deserve monitoring wherever expensive data are reused.

Limitations bound these claims: the transfer results are single-seed, five-epoch measurements at small scale, the RL comparison is a static verifier-selected proxy rather than on-policy policy optimization, and the DeepSeek-style replication tests components jointly rather than attributing the effect to any one of them. Whether response-only instruction tuning or production-scale mixtures damage specific alternative continuations remains open. Methodologically, the workflow used here---anomaly, aligned internal observables, targeted intervention, minimal reconstruction---turns a search-selected modification into a testable scientific question, and underlines that a measured benefit need not transfer across data regimes, architectures, or training stages.
